\documentclass[11pt,a4paper]{article}

\usepackage[T1]{fontenc}
\usepackage[utf8]{inputenc}
\usepackage[top=0.6in,bottom=0.9in,left=1in,right=1in]{geometry}
\usepackage{graphicx}
\usepackage{booktabs}
\usepackage{hyperref}
\usepackage{xcolor}
\usepackage{fancyvrb}
\usepackage{enumitem}
\usepackage{parskip}
\usepackage{titlesec}

\titlespacing*{\section}{0pt}{10pt}{4pt}
\titlespacing*{\subsection}{0pt}{7pt}{2pt}

\hypersetup{colorlinks=true, linkcolor=blue!60!black, urlcolor=blue!60!black}

\DefineVerbatimEnvironment{codeblock}{Verbatim}{
    fontsize=\footnotesize, frame=single, framesep=3pt, rulecolor=\color{gray!60}
}

\title{\vspace{-1.2cm}\textbf{INDENG 231 Course Project 2}\\[0.2em]
\large Using LLMs to Build a Personal Knowledge Base: Food Delivery Dispatch Algorithms}
\author{Anna Chen}
\date{May 2026}

\begin{document}
\maketitle

% ─────────────────────────────────────────────────────────────────────────────
\section{What I Built}
% ─────────────────────────────────────────────────────────────────────────────

I built a lightweight LLM-powered knowledge base system focused on the topic
of \textbf{food delivery platform dispatch algorithms} --- the real-time
decision-making pipelines that platforms like Meituan, Uber Eats, and DoorDash
use to match orders to riders, optimize delivery routes, and pre-position idle
couriers before demand spikes.
I chose this topic because it sits at an interesting intersection of combinatorial
optimization, machine learning, and real-time systems engineering, and because
it connects directly to the themes in Project 1.

\subsection{System Architecture}

The system follows a three-stage pipeline inspired by Karpathy's workflow:

\begin{codeblock}
raw/  -->  ingest.py (LLM)  -->  wiki/INDEX.md
                                  wiki/concepts/   [[backlinked articles]]
                                  wiki/algorithms/
                              -->  qa.py (Q&A CLI)
\end{codeblock}

\paragraph{Ingestion (\texttt{ingest.py}).}
The script reads each file in \texttt{raw/} and sends it to Claude (claude-opus-4-7)
with a prompt asking it to produce a clean wiki article: structured with headings,
a one-sentence \texttt{summary:} tag, and \texttt{[[concept]]} backlinks to related topics.
The LLM output is saved directly into the appropriate \texttt{wiki/} subdirectory and
a one-line entry is appended to \texttt{INDEX.md}.
Because the prompt is the same for every source, adding new articles requires no
changes to the pipeline --- just drop a new file into \texttt{raw/} and run
\texttt{python ingest.py}.

\paragraph{Q\&A (\texttt{qa.py}).}
A simple CLI that answers natural-language questions against the wiki.
When an \texttt{ANTHROPIC\_API\_KEY} is available, it selects the top-4 most relevant
articles (by keyword overlap) and sends them together with the question to Claude for a
synthesised answer.
Without an API key, it falls back to displaying relevant article excerpts directly ---
making the system usable even without spending tokens on every query.

To try it, run \texttt{python qa.py} for interactive mode, or pass a question directly:

\begin{codeblock}
python qa.py "How does batch matching work and why is it better than greedy?"
python qa.py "How does reinforcement learning apply to dispatch?"
python qa.py "How do platforms pre-position idle riders before demand spikes?"
\end{codeblock}

\subsection{Knowledge Base Content}

I seeded the \texttt{raw/} directory with four hand-written source documents
summarising the academic and industry literature on dispatch:

\begin{enumerate}[leftmargin=*, itemsep=1pt]
    \item \textbf{Dispatch Overview}: the three sub-problems (matching, routing, zone management),
          key metrics, and industry landscape.
    \item \textbf{Matching Algorithms}: greedy, Hungarian algorithm, auction algorithm,
          batch matching, and integer programming formulations.
    \item \textbf{Routing \& Batching}: VRP variants, order batching trade-offs,
          2-opt heuristics, dynamic re-routing, ETA prediction.
    \item \textbf{ML \& Reinforcement Learning}: demand forecasting, RL-based dispatch
          (DiDi KDD 2018, Meituan), GNNs for road graphs, fairness concerns.
\end{enumerate}

The LLM compiled these into four interlinked wiki articles totalling around 1,800 words,
automatically adding backlinks like \texttt{[[Bipartite Matching]]} and
\texttt{[[ETA Prediction]]} that connect related concepts across articles.

% ─────────────────────────────────────────────────────────────────────────────
\section{What Worked}
% ─────────────────────────────────────────────────────────────────────────────

\paragraph{LLM-generated structure was immediately useful.}
The biggest surprise was how much the LLM improved the \emph{organisation} of the
raw material, not just its readability.
My raw sources were dense and unstructured; Claude consistently identified which ideas
deserved their own heading, surfaced comparisons I hadn't made explicit (e.g., online
vs. offline matching, greedy vs. batch), and added cross-links that helped me see
connections I'd missed.

\paragraph{Backlinks created a genuinely navigable graph.}
The \texttt{[[concept]]} convention is simple but powerful.
After ingestion, ETA Prediction appeared as a backlinked concept in three separate articles,
which made it immediately obvious that this is a central dependency of the whole dispatch
system --- not just a detail in the routing section.
In Obsidian, these render as clickable graph links that make the knowledge base feel
like a real wiki rather than a pile of notes.

\paragraph{The keyword-based Q\&A fallback is more useful than expected.}
Without an API key, \texttt{qa.py} returns the top-2 relevant article excerpts.
For targeted questions (``how does batch matching work?''), this was often enough to
get a complete answer, because the wiki articles are already well-structured.
The LLM-powered mode would synthesise a tighter, more direct answer --- but the
fallback mode proves the value of having a well-organised wiki independently of the Q\&A layer.

% ─────────────────────────────────────────────────────────────────────────────
\section{What I Would Improve}
% ─────────────────────────────────────────────────────────────────────────────

\paragraph{Automatic source collection.}
Right now, source documents are written by hand.
A more complete system would include a scraper or web-clipper that pulls
arXiv abstracts, blog posts, or documentation pages directly into \texttt{raw/},
automatically converts them to clean markdown (e.g., using \texttt{trafilatura}
or the Obsidian Web Clipper), and triggers \texttt{ingest.py} to update the wiki.

\paragraph{Incremental updates instead of full rewrites.}
The current \texttt{ingest.py} creates a new wiki article from scratch each time.
A better approach would be to compare the new source against the existing wiki page
and only update the sections that changed --- reducing token usage and preserving
manual annotations.

\paragraph{Richer Q\&A with citation.}
The current Q\&A prompt just asks Claude to answer from the wiki.
A more robust version would ask it to cite specific article sections in its answer,
making it easy to verify and trace claims back to sources.
Retrieval-Augmented Generation (RAG) with embedding-based retrieval (instead of
keyword overlap) would also improve relevance on longer, more nuanced questions.

\paragraph{LLM ``health checks'' on the wiki.}
Following Karpathy's suggestion, I would add a \texttt{lint.py} script that
periodically runs Claude over the entire wiki looking for inconsistencies
(e.g., two articles giving conflicting definitions of the same concept),
orphaned backlinks (a \texttt{[[concept]]} that has no corresponding article),
and gaps worth filling (concepts referenced but never elaborated).

\paragraph{Visual output formats.}
Rendering outputs as Marp slides or matplotlib diagrams (e.g., a flowchart of the
dispatch pipeline) and filing them back into the wiki would make the knowledge base
much more useful for presentations and further study.

\end{document}
